\subsection{Das Konzept der Leichten Sprache und verfügbare Datenbasis}
\label{sec:konzept_leichte_sprache}

Das Konzept der Leichten Sprache (LS) stellt eine stark reglementierte Sprachvariante der deutschen Standardsprache dar, die primär darauf abzielt, Barrierefreiheit in der schriftlichen Kommunikation für Menschen mit kognitiven Einschränkungen, Lernschwierigkeiten oder geringer Lese- und Schreibkompetenz zu gewährleisten. Syntaktisch zeichnet sich LS unter anderem durch kurze Hauptsätze (Subjekt-Prädikat-Objekt), den Verzicht auf Nebensätze, Passivkonstruktionen, den Genitiv sowie den Konjunktiv aus. Auf lexikalischer Ebene werden abstrakte Begriffe durch konkrete Beispiele ersetzt, und lange zusammengesetzte Wörter werden mittels typografischer Hilfsmittel wie Bindestrichen oder dem Mediopunkt (z.~B. \textit{Bundestags-Wahl} oder \textit{Bundestags·wahl}) segmentiert. 

Vom Konzept der Leichten Sprache ist die sogenannte \textit{Einfache Sprache} (Plain German) abzugrenzen. Während die Leichte Sprache als hochspezialisierte Sprachvariante mit strengen, normierten Richtlinien konzipiert ist und in Deutschland rechtlich durch die Barrierefreie-Informationstechnik-Verordnung (BITV 2.0) für öffentliche Stellen vorgeschrieben wird \citep{bund2026barrierefreiheit}, besitzt die Einfache Sprache ein weitaus weniger restriktives Regelwerk. Letztere richtet sich an eine breitere Zielgruppe (wie Menschen mit Deutsch als Zweitsprache oder funktionale Alphabeten) und lässt komplexere Satzstrukturen sowie einen größeren Wortschatz zu. Ein weiterer kritischer Unterschied liegt im Qualitätssicherungsprozess: Texte in Leichter Sprache müssen zwingend von Personen aus der primären Zielgruppe auf ihre Verständlichkeit hin geprüft werden, was bei Texten in Einfacher Sprache nicht vorgeschrieben ist. Einen guten Überblick über den aktuellen Stand und die verfügbaren Ressourcen für Leichte und Einfache Sprache in Deutschland geben \citet{madina2023easy}. Sie zeigen dabei vor allem auf, wie unterschiedlich die bestehenden Datensätze strukturiert sind und dass es an standardisierten Werkzeugen zur automatischen Bewertung der Texte fehlt.

Die Erforschung datengetriebener, maschineller Übersetzungsverfahren (Automatic Text Simplification, ATS) für das Deutsche litt historisch unter einer eklatanten Datenknappheit. Erste Versuche, ein paralleles Korpus aufzubauen, gehen auf \citeauthor{klaper2013klaper} (\citeyear{klaper2013klaper}) mit dem \textit{KLAPER}-Korpus zurück, das jedoch mit ca. 250 parallelen Textpaaren (rund 7.000 Sätze und 70.000 Tokens) für moderne Deep-Learning-Ansätze deutlich zu klein war. 

Einen Meilenstein stellte die Arbeit von \citet{battisti2020corpus} dar, die ein umfangreiches, multimodales Korpus für das Deutsche zusammentrugen. Die AutorInnen entwickelten Web-Scraper, um Texte in Leichter und Einfacher Sprache von verschiedenen deutschsprachigen Webportalen und Institutionen der Behindertenhilfe (wie etwa gesundheitlichen Informationsdiensten, öffentlich-rechtlichen Nachrichtenportalen und sozialen Trägern) zu extrahieren. Insgesamt umfasst das Korpus 6.217 Dokumente bzw. fast 211.000 Sätze mit rund 5,5 Millionen geschriebenen Tokens. Ein zentrales Merkmal der Arbeit ist die Aufteilung des Korpus: Es besteht zu einem großen Teil aus 5.461 monolingualen Dokumenten in vereinfachter Sprache, während der parallele Teil mit lediglich 378 parallelen Dokumentenpaaren (jeweils 378 Ausgangs- und Zieltexte) im Vergleich dazu sehr klein ausfällt. Neben den reinen Texten erfassten die AutorInnen systematisch typografische Metadaten (Schriftart, -größe und -stil), Absatz- und Zeilenstrukturen sowie Informationen zu eingebetteten Bildern (Vorhandensein, Dimensionen und Positionierung). Sie wiesen nach, dass diese visuellen Layout-Features die automatische Bewertung der Lesbarkeit (Readability Assessment) stark verbessern. 

Für den Bereich der "`Einfachen Sprache"' (Plain German) präsentierten \citet{stodden2023deplain} mit \textit{DEplain} ein professionell erstelltes paralleles Korpus. Dieses teilt sich auf in einen Nachrichtenbereich (DEplain-APA) mit ca. 500 Dokumentenpaaren (rund 13.000 manuell ausgerichteten Satzpaaren) und einen Webbereich mit ca. 150 manuell ausgerichteten Dokumenten (ca. 2.000 Satzpaaren). Um die Datenknappheit weiter zu reduzieren, stellt das Framework zusätzlich einen Web-Harvester zur automatischen Extraktion und Ausrichtung bereit, wodurch der Webbereich auf ca. 750 Dokumentenpaare (ca. 3.500 automatisch ausgerichtete Satzpaare) vergrößert werden kann. Die Autoren zeigten, dass Modelle, die auf Dokumentenebene trainiert werden, oft Probleme haben, die feinkörnigen semantischen Übergänge beizubehalten, weshalb präzise Alignments eine kritische Ressource für das Training von Seq2Seq-Modellen darstellen. An dieser Stelle ist jedoch hervorzuheben, dass eine direkte Vergleichbarkeit mit Systemen zur Übersetzung in Leichte Sprache erschwert ist: Obwohl das DEplain-Framework methodisch wertvolle Ansätze liefert, bezieht es sich durch den Fokus auf Einfache Sprache auf eine andere linguistische Ziel- und Datendomäne.

Einen weiteren Beitrag zur Generierung paralleler Ressourcen leisteten \citet{toborek2023new} mit der Vorstellung des \textit{New Aligned Simple German Corpus}. Ihr Datensatz besteht aus ca. 700 Dokumentenpaaren auf Dokumentenebene und über 10.000 Satzpaaren auf Satzebene, die teils manuell annotiert und teils automatisch berechnet wurden. Das Korpus kombiniert dabei das systematische Scraping paralleler Webinhalte mit automatischer Satz-Ausrichtung. Allerdings fasst die Arbeit Texte sowohl der Leichten als auch der Einfachen Sprache unter einem gemeinsamen Vereinfachungsbegriff zusammen, was eine feine Unterscheidung der spezifischen syntaktischen und semantischen Richtlinien von Leichter Sprache erschwert.

\subsection{Maschinelle Textvereinfachung und Satz-Alignment}
\label{sec:maschinelle_textvereinfachung_alignment}

Die maschinelle Textvereinfachung lässt sich formal als Monolingual Machine Translation begreifen, bei der ein komplexer Ausgangstext in eine vereinfachte Zielsprachvariante übersetzt wird. Eine fundamentale Herausforderung im Vergleich zur klassischen maschinellen Übersetzung (z.~B. Deutsch-Englisch) ist die Asymmetrie der Satzstrukturen: Eine Vereinfachung erfordert häufig Satzaufspaltungen (Sentence Splitting), Auslassungen irrelevanter Informationen (Compression) oder Paraphrasierungen. Folglich liegt selten ein einfaches 1:1-Satzalignment vor; stattdessen dominieren 1:n-, n:1- oder n:m-Alignments.

Zur automatisierten Ausrichtung paralleler Dokumente auf Satz- und Absatzebene hat sich das Werkzeug \textit{CATS} von \citet{stajner2018cats} etabliert. CATS implementiert verschiedene lexikalische und semantische Ähnlichkeitsmetriken (wie Character N-Grams (CNG) und gemittelte Word-Embeddings) und kombiniert diese mit dem \textit{Most Similar Text} (MST) bzw. \textit{MST-LIS} (Longest Increasing Subsequence) Suchalgorithmus. \citet{stajner2018cats} wiesen nach, dass die Einbeziehung semantischer Vektoren (WAVG) und kontinuierlicher Wort-Alignments (CWASA) im Vergleich zu rein lexikalischen Maßen robustere Alignments bei stark paraphrasierten Sätzen liefert, wie sie für Leichte Sprache typisch sind. Dennoch bleibt die automatische Filterung verrauschter Alignments auf Basis kosinusähnlicher Schwellenwerte ein aktives Forschungsfeld, da zu niedrige Schwellenwerte fehlerhafte Übersetzungen und zu hohe Schwellenwerte einen Informationsverlust im Trainingsdatensatz induzieren.

\subsection{Evaluierungsmetriken für die Vereinfachung}
\label{sec:evaluierungsmetriken}

Die automatische Evaluierung von Textvereinfachungssystemen stellt eine große methodische Hürde dar. Traditionelle Metriken der maschinellen Übersetzung wie BLEU, die primär auf n-Gramm-Präzision im Vergleich zu einer Referenzübersetzung basieren, erweisen sich für die Vereinfachung oft als ungeeignet. Wie \citet{alva2021unsuitability} empirisch belegen, korreliert BLEU nur schwach mit dem menschlichen Urteil über Einfachheit und Textqualität. Der Grund liegt darin, dass BLEU lexikalische Variationen und strukturelle Modifikationen (wie Satzteilung und Umformulierung) als Fehler abstraft, obwohl genau diese für eine gute Vereinfachung essenziell sind.

Als Alternative etablierte sich die Metrik \textit{SARI} (System Accessibility and Readability Index), eingeführt von \citet{xu2016optimizing}. SARI vergleicht die Systemausgabe sowohl mit dem komplexen Ausgangstext als auch mit den Referenzübersetzungen. Dabei berechnet SARI separat die Präzision und den Recall für eingefügte (add), gelöschte (delete) und beibehaltene (keep) Wörter:
\begin{equation}
    \text{SARI} = \frac{F_{\text{add}} + F_{\text{keep}} + P_{\text{del}}}{3}
\end{equation}
Dadurch belohnt SARI explizit Vereinfachungsoperationen, die sich vom komplexen Ausgangstext wegbewegen.

Für die deutsche Sprache stellt die neu entwickelte Metrik \textit{DETECT} von \citet{stodden2026detect} einen wichtigen Fortschritt dar. DETECT versucht, die Dimensionen Verständlichkeit (Ease), Grammatikalität und semantische Konsistenz (Textual Clarity) über ein erlerntes Evaluierungsmodell abzubilden, das auf synthetischen LLM-Generaten trainiert wurde. Dies adressiert die Schwäche klassischer Referenzmetriken, die die semantische Abweichung (Information Loss) bei extremen Kürzungen oft nicht hinreichend quantifizieren können.

\subsection{Präferenzoptimierung (DPO) in der Textgenerierung}
\label{sec:dpo_textgenerierung}

Während Supervised Fine-Tuning (SFT) darauf abzielt, die Token-Verteilung eines Referenzkorpus zu imitieren, stößt dieser Ansatz bei der Durchsetzung spezifischer Stil- und Formatierungsregeln (wie den Richtlinien der Leichten Sprache) an Grenzen. Um generative Modelle präziser auf menschliche Qualitätsurteile und formale Textregeln auszurichten, wurde traditionell Reinforcement Learning from Human Feedback (RLHF) über Proximal Policy Optimization (PPO) eingesetzt. Dieser Prozess ist jedoch rechenzeitintensiv und hyperparameter-instabil.

Mit der Einführung von \textit{Direct Preference Optimization} (DPO) zeigten \citet{rafailov2023direct}, dass die explizite Modellierung eines separaten Reward-Modells mathematisch umgangen werden kann. DPO optimiert die Policy direkt auf Präferenzdaten der Form $(x, y_w, y_l)$ — bestehend aus einem Prompt $x$, einer bevorzugten Antwort $y_w$ (winning) und einer abgelehnten Antwort $y_l$ (losing) — mittels einer geschlossenen Verlustfunktion:
\begin{equation}
    \mathcal{L}_{\text{DPO}}(\theta; \pi_{\text{ref}}) = - \mathbb{E}_{(x, y_w, y_l)} \left[ \log \sigma \left( \beta \log \frac{\pi_\theta(y_w | x)}{\pi_{\text{ref}}(y_w | x)} - \beta \log \frac{\pi_\theta(y_l | x)}{\pi_{\text{ref}}(y_l | x)} \right) \right]
\end{equation}
wobei $\pi_{\text{ref}}$ die SFT-Referenz-Policy darstellt, $\beta$ ein Regularisierungsparameter ist und $\sigma$ die Sigmoidfunktion bezeichnet.

Die Anwendung von DPO auf deutsche Textvereinfachung wurde kürzlich von \citet{dpo2025personalizing} untersucht. Die Autoren nutzten den \textit{HF4ATS}-Datensatz (Human Feedback for Automatic Text Simplification), um LLMs gezielt auf die Präferenzen von Menschen mit kognitiven Einschränkungen auszurichten. Sie wiesen nach, dass DPO-trainierte Modelle im Vergleich zu reinem SFT nicht nur eine signifikant höhere subjektive Verständlichkeit erzielen, sondern auch formale Kriterien der Leichten Sprache (wie das Vermeiden von Genitiven und Nebensätzen) stabiler einhalten, ohne dabei halluzinierte Fakten oder unzusammenhängende Sätze zu erzeugen. Dies motiviert den in dieser Arbeit gewählten Ansatz, DPO unter Verwendung eines trainierten Komplexitätsklassifikators als automatische Reward-Instanz zur Optimierung einzusetzen.
